\subsection{The mRNA H-candidate projection for the $B^*_{Q6}$ abundance residual}\label{subsec:bior-hcandidate-mrna-bstar-q6-abundance-residual}\origin{ai}

\paragraph{Finite panel.} The record studies a finite panel of coding sequences in which nine synonymous residual coordinates are read against protein abundance, modeled translation weights, and matched mRNA abundance.  The coordinate list is
$K_{AAA}$, $Arg_{AGR}$, $Ile_{AUA}$, $Leu_{CUN\,vs\,UUR}$, $Leu_{UUA\,vs\,UUG}$, $Ser_{UCR\,vs\,AGY}$, $Ser_{UCA\,vs\,UCG}$, $Thr_{ACR\,vs\,ACY}$, and $f_3$-stress.  The control block contains an intercept, $\log(\mathrm{cds\_len\_nt})$, the twenty amino-acid composition proportions $A,C,D,E,F,G,H,I,K,L,M,N,P,Q,R,S,T,V,W,Y$, $GC3$, and an $M$-density baseline defined as the fraction of sense codons on the union support of the nine projected $B^*_{Q6}$ vectors.  In each organism the declared control column count is $24$ and the observed control rank is $23$, so the projection is not an unconstrained restatement of the response column.

\paragraph{Joined rows and rank.} The finite row counts are organism-specific.  Danio rerio has $n=13140$ rows after the mRNA join, with $14410$ CDS/protein-abundance rows available before requiring mRNA and $\operatorname{rank}(Z)=23$.  Escherichia coli K--12 MG1655 has $n=3415$, with $3739$ CDS/protein rows, $3747$ valid abundance records, and $\operatorname{rank}(Z)=23$.  Homo sapiens has $n=11554$, with $19053$ CDS/protein rows, $19483$ valid abundance records, and $\operatorname{rank}(Z)=23$.  Saccharomyces cerevisiae has $n=5128$, with $5815$ CDS/protein rows, $6351$ valid abundance records, and $\operatorname{rank}(Z)=23$.  Thus the mRNA projection is a row-restricted computation: it asks what part of the modeled-translation complement remains aligned with transcript abundance after the matched mRNA rows are imposed, not what would be inferred on the larger protein-only join.

\paragraph{Modeled complement before mRNA.} Let $P_Q$ denote the abundance residual component associated with the $B^*_{Q6}$ coordinate block after the stated controls, and let $P_{Q\perp T_{mod}}$ denote the part left after the modeled translation-weight component is removed.  The recorded complement ratios are
$$\begin{aligned}
 d_{\mathrm{Danio}} &= 0.8424312412078521,\\
 d_{\mathrm{E.coli}} &= 0.6628211177352172,\\
 d_{\mathrm{Human}} &= 0.8566190269974241,\\
 d_{\mathrm{Yeast}} &= 0.22386566043976477.
\end{aligned}$$
The corresponding squared norms are concrete: for Danio rerio, $\|P_Q\|^2=1092.0114815505467$ and $\|P_{Q\perp T_{mod}}\|^2=919.9445878158525$; for E. coli, $393.78692095732885$ and $261.0102870984463$; for human, $565.3850201235074$ and $484.31956581711796$; for yeast, $675.3275479192903$ and $151.18264752811882$.  Yeast is therefore the strongest modeled-translation absorption case in this finite panel, while Danio and human retain a large complement after the modeled readout.

\paragraph{mRNA H-candidate readback.} The mRNA candidate is then tested by projecting $P_{Q\perp T_{mod}}$ onto residualized $\log_{10}$ mRNA under the same control discipline.  The measured mRNA absorption values are small in all four organisms:
$$\begin{aligned}
 a_{\mathrm{Danio}} &= 0.011973033879985792, & d_{\mathrm{combined}} &= 0.8322280386743373,\\
 a_{\mathrm{E.coli}} &= 0.0017706954554908886, & d_{\mathrm{combined}} &= 0.6615987644523418,\\
 a_{\mathrm{Human}} &= 0.007858222847072218, & d_{\mathrm{combined}} &= 0.8498079619223776,\\
 a_{\mathrm{Yeast}} &= 0.0018976054721605224, & d_{\mathrm{combined}} &= 0.223360891719316.
\end{aligned}$$
These values make the H-candidate result a descriptive partition of variance, not a causal transcript mechanism.  The largest mRNA absorption is in Danio rerio, about $1.20\%$ of the modeled-translation complement, and the other three organisms are below $0.79\%$.  In particular, the finite data do not support the claim that transcript abundance explains the residual complement; they only certify that the mRNA projection has been computed and bounded on the joined rows.

\paragraph{Join self-check.} The mRNA restriction changes the row universe in a visible way.  Comparing the current all-readout join to the protein-abundance complement reference, E. coli differs by only $0.00017488612263594217$ and yeast by $0.012133464712697745$, both within the $0.02$ tolerance.  Danio differs by $0.03219747984141397$ and human by $0.06427590857052146$, outside that tolerance.  This is not a biological failure signal; it is a finite-row warning that mRNA availability changes the protein-id join in those two organisms.  The stated readback is therefore row-conditioned: the mRNA H-candidate numbers belong to the mRNA-matched panels, not to the larger protein-only panels.

\paragraph{Relation to half-life mediation rows.} A separate cross-domain computation supplies a useful contrast.  In an E. coli half-life panel with $2720$ joined stable identifiers, the reported CDS hit rate is $0.9978649586335735$, the joint $B^*_{Q6}$ mediation fraction is $0.24697202727892137$, and the right-tail permutation value is $0.000999000999000999$.  In the human half-life panel with $16165$ joined stable identifiers and CDS hit rate $0.9779294769799313$, the corresponding mediation fraction is $0.1411299863319291$, with the same recorded right-tail permutation value $0.000999000999000999$.  Those half-life rows concern an RNA-stability readout and are not the same object as the mRNA-abundance H-candidate projection here.  They show that RNA-layer decompositions can be finite and nontrivial, while the present mRNA-abundance projection remains small under the declared controls.

\paragraph{Local verdict.} The finite claim is ready as a bounded cross-layer relation: $B^*_{Q6}$ abundance residuals can be decomposed into a modeled-translation complement and a further mRNA-abundance projection on organism-specific joined panels.  The numerical result is negative for strong mRNA absorption, because the computed transcript component is small in every organism and the row-join self-check fails tolerance in Danio and human.  No translation survival, folding fact, physical admissibility fact, molecular function, mechanism, phenotype, or universal biological rule follows from this chapter.  Such a promotion would require a separate layer-matched contact, such as a direct translation-efficiency survival matrix $S^{QT}$ or a controlled protein/readout mediation matrix $M^{QTP}$, computed on the relevant finite rows.